



\Large{\textbf{Status of the work from $25^{th}$ October 2021 till date}}
\begin{enumerate}
    \item \large{\textbf{Completed Course Work}}
    \begin{itemize}
        \item Research Methods in Computer Science - Core Course 
        \item Data Structures And Programming Lab - Core Course
        \item Algorithms - Core Course
        \item Research Publications and Ethics - Core Course
        \item Machine Learning - Elective Course
        \item Social Network Analysis - Audit Course
    \end{itemize}
    \item \large{\textbf{Survey of Visualization Recommendation and NLDR approaches}}
    \begin{itemize}
        \item Extensive survey on Visualization recommendation \sloppy{approaches} can be categorized into visualization selection, technique, and recommendation.
        \item Extension survey on nonlinear dimensionality reduction approaches which can visualize the high-dimensional data.
    \end{itemize}
    
    \item \large{\textbf{Contribution 1:}}\\
    \justifying{Kuruva, H., Swarupa Rani, K., Moiz, S.A. (2023). \textbf{Visualization Recommendation for Incremental Data Based on Intent.} In: Morusupalli, R., Dandibhotla, T.S., Atluri, V.V., Windridge, D., Lingras, P., Komati, V.R. (eds) Multi-disciplinary Trends in Artificial Intelligence. MIWAI 2023. Lecture Notes in Computer Science, vol 14078. Springer, Cham. \url{https : //doi.org /10.1007/978−3−031−36402−0 26}}
    
    \item \large{\textbf{Contribution 2:}}\\
    Proposed $HM\_QCv$: A Computationally Efficient Hybrid Model for Enhancing Quality Clusters in Visualization - \textbf{Communicated to Big Data Research Journal (Special Issue) on $21^{st}$ Jan 2025.}
    \item \large{\textbf{Contribution 3:}}\\
    Proposing a framework for Incremental Visualization Recommendation using Feature Selection.
\end{enumerate}


The contributions are discussed in the Sections 3, 4 and 5.